% Unofficial Umich Poster template.
% A fork of the MSU template https://www.overleaf.com/latex/templates/an-unofficial-poster-template-for-michigan-state-university/wnymbgpxnnwd
% which is a fork of https://www.overleaf.com/latex/templates/an-unofficial-poster-template-for-new-york-university/krgqtqmzdqhg
% which is a fork of https://github.com/anishathalye/gemini
% also refer to https://github.com/k4rtik/uchicago-poster



\documentclass[final]{beamer}

% ====================
% Packages
% ====================

\usepackage[T1]{fontenc}
 \usepackage[utf8]{luainputenc}
\usepackage{lmodern}
% UK A1 landscape dimensions are 84.1cm x 59.4cm
\usepackage[size=custom, width=84,height=59.4, scale=1.2]{beamerposter}
\usetheme{gemini}
\usecolortheme{msu}
\usepackage{graphicx}
\usepackage{booktabs}
\usepackage{tikz}
\usepackage{pgfplots}
\pgfplotsset{compat=1.14}
\usepackage{anyfontsize}

% ====================
% Lengths
% ====================

% If you have N columns, choose \sepwidth and \colwidth such that
% (N+1)*\sepwidth + N*\colwidth = \paperwidth
\newlength{\sepwidth}
\newlength{\colwidth}
\setlength{\sepwidth}{0.025\paperwidth}
\setlength{\colwidth}{0.3\paperwidth}

\newcommand{\separatorcolumn}{\begin{column}{\sepwidth}\end{column}}

% ====================
% Title
% ====================

\title{Measuring Competition and Cooperation in LLM Bargaining: An Empirical Meta-Game Analysis}

\author{Gabriel Smithline, Chris Mascioli, Mithun Chakraborty, and Michael P. Wellman}

\institute[shortinst]{School of Information, University of Michigan}

% ====================
% Footer (optional)
% ====================

\footercontent{
  \href{https://github.com}{Github: https://github.com} \hfill
Poster Session \hfill
  \href{mailto:youremail@msu.edu}{youremail@umich.edu}}
% (can be left out to remove footer)

% ====================
% Logo (optional)
% ====================

% use this to include logos on the left and/or right side of the header:
% Left: institution
 \logoright{
\includegraphics[height=5cm]{logos/blockM.png}}
% Right: funding agencies and other affilations 
%\logoright{\includegraphics[height=7cm]{logos/NSF.eps}}
% ====================
% Body
% ====================

\begin{document}



\begin{frame}[t]
\begin{columns}[t]
\separatorcolumn

\begin{column}{\colwidth}

  \begin{block}{Motivation}
  \begin{itemize}
    \item Performance of an LLM agent hinges on the behavior of its counterpart.
    \item Sequential bargaining offers a mixed-motive testbed with measurable
          welfare and fairness.
    \item We study how prompting and simple heuristics shape negotiation
          outcomes.
  \end{itemize}
  \end{block}

  \begin{block}{Background}
  We evaluate bargaining performance of large language models (LLMs) using
  empirical meta-game analysis. Agents negotiate over five items with stochastic
  outside options in three games:\@\,BG1 ($\gamma=0.90,R=3$), BG2
  ($\gamma=0.98,R=3$), and BG3 ($\gamma=0.98,R=5$). From 1,000 bootstrap
  resamples per game we compute a maximum-entropy Nash equilibrium (MENE) to
  benchmark regret, welfare, and EF1 fairness.
  \end{block}

  \begin{block}{Key Contributions}
  \begin{itemize}
    \item Empirical meta-game built from 35k+ simulated bargaining games.
    \item Seven-level prompting hierarchy reveals how guidance shapes strategy.
    \item Multi-criteria analysis connects regret, welfare, and fairness.
    \item Robustness checks expose vulnerabilities to rigid heuristics.
  \end{itemize}
  \end{block}

  \begin{block}{Experimental Setup}
  \begin{itemize}
    \item 9 LLM strategies combining models from OpenAI, Anthropic, and Google.
    \item Heuristic baselines {\small (tough, soft)} and NFSP learning agent.
    \item 35k+ simulated games \textrightarrow{} empirical meta-games and BG0
          (mixture over BG1--BG3).
  \end{itemize}
  \end{block}

\end{column}

\separatorcolumn

\begin{column}{\colwidth}

  \begin{block}{LLM--Only Results}
  Table~1 reports bootstrap means for regret, Nash and utilitarian welfare, and
  EF1 frequency across BG1--BG3 and BG0. OpenAI\textapos{}s
  \texttt{o3-mini-c-0} consistently achieves the lowest regret (2--5) and high
  welfare ($\approx$92\% NW) while Gemini models show higher regret and lower
  fairness. GPT--4o variants form the next best group.
  \end{block}

  \begin{figure}
    \centering
    \includegraphics[width=\linewidth]{icml_2025_workshop/average_br_graphs/llm_only/avg_br_graph_gm_full_llm_only.png}
    \caption{BG0 best-response graph for LLMs.}
  \end{figure}

  \begin{figure}
    \centering
    \includegraphics[width=0.48\linewidth]{icml_2025_workshop/circles_dominiated_mistakes/4o_dominated_mistakes.png}
    \hfill
    \includegraphics[width=0.48\linewidth]{icml_2025_workshop/circles_dominiated_mistakes/sonnet_3.7_dominated_mistakes.png}
    \caption{Dominated circle mistake rates for GPT--4o (left) and Sonnet~3.7 (right).}
  \end{figure}

\end{column}

\separatorcolumn

\begin{column}{\colwidth}

  \begin{block}{Mixed Agents}
  Table~2 adds heuristic and NFSP baselines. \texttt{tough} and \texttt{soft}
  are easily exploited, showing regret $>80$ and near-zero EF1. LLMs still
  dominate; \texttt{o3-mini-c-0} and GPT--4o variants retain low regret and high
  welfare in BG0.
  \end{block}

  \begin{figure}
    \centering
    \includegraphics[width=\linewidth]{icml_2025_workshop/average_br_graphs/llm_+_non_llm/avg_br_graph_gm_4_llm_and_non_llm_b.png}
    \caption{BG0 best-response graph with LLM and non-LLM agents.}
  \end{figure}

  \begin{block}{Conclusion}
  \begin{itemize}
    \item OpenAI models provide strong bargaining strategies with low regret and
          high fairness.
    \item Prompt engineering (circles) significantly impacts performance of
          Anthropic and Google models.
    \item Heuristic baselines underperform dramatically.
  \end{itemize}
  \end{block}

\end{column}

\separatorcolumn
\end{columns}
\end{frame}

\end{document}
